The condition is equivalent to all subset sums being distinct: if two subsets share a sum, remove their intersection to obtain disjoint subsets with equal sums. Build the possible sets in increasing order, keeping a bitset $B$ of attainable sums, initially only bit 0. To append $a$, the old and new sums are disjoint exactly when $B\mathbin{\&}(B\ll a)=0$; update $B$ to $B\mathbin{|}(B\ll a)$. Exhaustively branching on inclusion or exclusion of each integer 1 through 15 yields the largest sums by set size $0,1,2,3,4,5$ as $0,15,29,42,53,61$; there are no valid larger sets. Thus the maximum is $\boxed{61}$, attained by $\{8,11,13,14,15\}$. The algorithm has only $2^{15}$ candidate sets, and the bitset disjointness test is exactly the proved condition, so it is an exhaustive certificate rather than a heuristic.
